% Template for Elsevier
% version 1.1-5p dated 18 January 2011
% NO ACCENTUATION PROBLEM UTF-8 : NOT ACCEPTED

% This file (c) 2010-2011 Elsevier Ltd. Modifications may be freely made,
% provided the edited file is saved under a different name

% This file contains modifications for Procedia Computer Science
% but may easily be adapted to other journals


%-----------------------------------------------------------------------------------

%% This template uses the elsarticle.cls document class and the extension package ecrc.sty
%% For full documentation on usage of elsarticle.cls, consult the documentation "elsdoc.pdf"
%% Further resources available at http://www.elsevier.com/latex

%-----------------------------------------------------------------------------------

%%%%%%%%%%%%%%%%%%%%%%%%%%%%%%%%%%%%%%%%%%%%%%
%%%%%%%%%%%%%%%%%%%%%%%%%%%%%%%%%%%%%%%%%%%%%%
%%                                          %%
%% Important note on usage                  %%
%% -----------------------                  %%
%% This file must be compiled with PDFLaTeX %%
%% Using standard LaTeX will not work!      %%
%%                                          %%
%%%%%%%%%%%%%%%%%%%%%%%%%%%%%%%%%%%%%%%%%%%%%%
%%%%%%%%%%%%%%%%%%%%%%%%%%%%%%%%%%%%%%%%%%%%%%

%% The '5p' and 'times' class options of elsarticle are used for Elsevier CRC
\documentclass[5p]{elsarticle}


%%%%%%%%%%%%%%%%%%%%%%%%%%%%%%%%%
\usepackage[spanish,es-noshorthands]{babel}     % Language
\usepackage[T1]{fontenc}
\usepackage{float}
\usepackage{parskip}
\usepackage[utf8]{inputenc}   % UTF-8 encoding
\usepackage{amsmath}             % For equation references with \eqref
\usepackage{flushend}            % To equalize columns on the last page
\usepackage[colorlinks=true,linkcolor=blue,citecolor=blue,urlcolor=blue]{hyperref}
\newcommand*{\doi}[1]{\href{http://doi.org/#1}{#1}}

%%%%%%%%%%%%%%%%%%%%%%%%%%%%%%%%%%%%%%%%%%%%%%%%%%%%%%

%% Hereafter the template follows `elsarticle'.
%% For more details see the existing template files elsarticle-template-harv.tex and elsarticle-template-num.tex.

%% Elsevier CRC generally uses a numbered reference style
%% For this, the conventions of elsarticle-template-num.tex should be followed (included below)
%% If using BibTeX, use the style file elsarticle-num.bst

%%%%%%%%%%%%%%%%%%%%%%%%%%%%%%%%%%%%%%%%%%%%%%%%%%%%%%%%%%%%%%%%%%%%%%%%%%

%% The amssymb package provides various useful mathematical symbols
\usepackage{amssymb}
%% The amsthm package provides extended theorem environments
%% \usepackage{amsthm}

%% The lineno packages adds line numbers. Start line numbering with
%% \begin{linenumbers}, end it with \end{linenumbers}. Or switch it on
%% for the whole article with \linenumbers after \end{frontmatter}.
%% \usepackage{lineno}

%% natbib.sty is loaded by default. However, natbib options can be
%% provided with \biboptions{...} command. Following options are
%% valid:

%%    round  -  round parentheses are used (default)
%%    square -  square brackets are used    [option]
%%    curly  -  curly braces are used      {option}
%%    angle  -  angle brackets are used    <option>
%%    semicolon  -  multiple citations separated by semi-colon
%%    colon  - same as semicolon, an earlier confusion
%%    comma  -  separated by comma
%%    numbers-  selects numerical citations
%%    super  -  numerical citations as superscripts
%%    sort   -  sorts multiple citations according to order in ref. list
%%    sort&compress    -  like sort, but also compresses numerical citations
%%    compress - compresses without sorting
%%
%% \biboptions{comma,round}


% if you have landscape tables
\usepackage[figuresright]{rotating}

% put your own definitions here:
%    \newcommand{\cZ}{\cal{Z}}
%    \newtheorem{def}{Definition}[section]
%    ...

% add words to TeX's hyphenation exception list

% to be able to introduce multiple figures that occupy the width of the two columns.
\usepackage{subfigure}
\usepackage{array}
\usepackage{placeins}

% declarations for front matter

\begin{document}

\tolerance=1000
\emergencystretch=3em

\begin{frontmatter}


%% Title, authors and addresses

%% use the tnoteref command within \title for footnotes;
%% use the tnotetext command for the associated footnote;

%% use the fnref command within \author or \address for footnotes;
%% use the fntext command for the associated footnote;

%% use the corref command within \author or corresponding author footnotes;
%% use the cortext command for the associated footnote;
%% use the ead command for the email address,
%% and the form \ead[url] for the home page:
%%
%% \title{Title\tnoteref{label1}}
%% \tnotetext[label1]{}
%% \author{Name\corref{cor1}\fnref{label2}}
%% \ead{email address}
%% \ead[url]{home page}
%% \fntext[label2]{}
%% \cortext[cor1]{}
%% \address{Address\fnref{label3}}
%% \fntext[label3]{}

%\dochead{Article Header}
%% Use \dochead if there is an article header, e.g. \dochead{Short communication}

\title{Sistema de Planificación de Viajes en Transporte Público con Validación Formal en VDM++ para la Ciudad de Arequipa}


%% use optional labels to link authors explicitly to addresses:
%% \author[label1,label2]{<author name>}
%% \address[label1]{<address>}
%% \address[label2]{<address>}

\author[First]{Integrante 1}

\author[First]{Integrante 2}

\author[First]{Integrante 3}

\author[First]{Integrante 4}

\author[First]{Integrante 5}

\address[First]{Universidad Nacional de San Agustín, Arequipa, Perú.}

\begin{abstract}
La información sobre las rutas de transporte público en Arequipa es dispersa e imprecisa, lo que provoca desplazamientos ineficientes y decisiones de movilidad erróneas. Este trabajo propone una plataforma web para la búsqueda de rutas directas y combinadas, cuya componente crítica, el algoritmo de planificación de viajes, es especificada y validada formalmente con el lenguaje VDM++. Se definen los tipos de datos e invariantes del dominio de rutas y se especifican las operaciones de búsqueda mediante pre y postcondiciones que garantizan propiedades como la distancia máxima de acceso a las rutas y la viabilidad de los transbordos. La propuesta se valida mediante la herramienta VDM++ Toolbox, abordando así tanto la correctitud del algoritmo como la escalabilidad de la solución.
\end{abstract}

\begin{keyword}
métodos formales \sep VDM++ \sep transporte público \sep planificación de rutas \sep sistemas críticos
\end{keyword}


%%
%% Start line numbering here if you want
%%
% \linenumbers
\end{frontmatter}


%% main text
\section{Introducción}
\label{sec:introduccion}

\subsection{Definición del Problema}
\label{sec:problema}

En la ciudad de Arequipa, la información sobre las rutas del transporte público es dispersa, está desactualizada y no se encuentra disponible de forma clara para el ciudadano. Ante la necesidad de desplazarse entre dos puntos, el usuario desconoce qué línea de bus tomar, cuáles son sus paraderos y si es necesario realizar un transbordo. Esta falta de información confiable produce pérdida de tiempo, desplazamientos innecesarios y decisiones de movilidad incorrectas.

La planificación de viajes en transporte público constituye una tarea computacionalmente \emph{crítica}: un resultado erróneo (una ruta que no pasa cerca del origen, un transbordo inviable o una dirección de recorrido incorrecta) induce al usuario a tomar una decisión equivocada, con el consiguiente costo en tiempo y seguridad. Por este motivo, la correctitud del algoritmo de búsqueda de rutas no puede garantizarse únicamente mediante pruebas de ejecución; se requiere una especificación formal que demuestre propiedades sobre el resultado producido, tal como se argumenta para los sistemas de implementación crítica~\cite{verhoef2005}.

El presente trabajo propone una plataforma web para la búsqueda de rutas de transporte en Arequipa. Su componente crítica, el algoritmo de búsqueda de rutas directas y combinadas, es especificada y validada formalmente mediante el lenguaje VDM++~\cite{fitzgerald2008}, utilizando invariantes, precondiciones y postcondiciones. La validación del modelo se realiza con la herramienta VDMTools~\cite{fitzgerald2008}.

\subsection{Objetivos}
\label{sec:objetivos}

\subsubsection{Objetivo general}

Desarrollar una propuesta de solución para la planificación de viajes en transporte público, cuya componente crítica, el algoritmo de búsqueda de rutas, sea especificada y validada formalmente con VDM++, garantizando su correctitud y escalabilidad.

\subsubsection{Objetivos específicos}

\begin{enumerate}
    \item Definir los tipos de datos e invariantes del dominio de rutas (recorridos, paraderos y radios de búsqueda).
    \item Especificar en VDM++ las operaciones de búsqueda de rutas directas y de rutas combinadas mediante pre y postcondiciones.
    \item Establecer las propiedades que el algoritmo debe garantizar, tales como la distancia máxima de acceso a una ruta, la viabilidad de los transbordos y la longitud mínima de los segmentos.
    \item Validar el modelo formal con la herramienta VDM++ Toolbox mediante obligaciones de prueba y casos de prueba.
    \item Documentar los requerimientos funcionales y los trabajos relacionados que sustentan la propuesta.
\end{enumerate}

\section{Trabajos Relacionados}
\label{sec:trabajos_relacionados}

La presente propuesta se fundamenta en la línea de investigación de la especificación formal con VDM++ y su aplicación a sistemas de implementación crítica y de gran escala.

Respecto al soporte de herramientas, Fitzgerald, Larsen y Sahara~\cite{fitzgerald2008} describen VDMTools, un conjunto de herramientas que analizan modelos expresados en VDM y sus dialectos VDM++ y VDM-RT, y que ofrecen verificación estática, generación automática de código, mapeo a UML, análisis de cobertura de pruebas y depuración. Este trabajo destaca que VDMTools está orientado al desarrollo rentable de modelos formales en entornos industriales, lo que sustenta la elección de esta herramienta para validar el algoritmo propuesto.

En el ámbito de los sistemas críticos, Verhoef, Larsen y Hooman~\cite{verhoef2006} extienden VDM++ para modelar y validar sistemas embebidos distribuidos en tiempo real, definiendo una semántica operacional constructiva validada con VDMTools. De forma complementaria, Verhoef~\cite{verhoef2005} analiza el uso de VDM++ para especificar sistemas de tiempo real, argumentando que los métodos formales deben integrarse a los flujos de trabajo industriales y permitir el análisis temprano de propiedades. Ambos trabajos justifican la aplicación de VDM++ a un componente cuya correctitud es crítica, como el algoritmo de planificación de viajes propuesto.

Finalmente, Mochio y Araki~\cite{mochio2011} vinculan VDM++ con el desarrollo formal ágil y escalable de software de misión crítica, incluyendo la descripción de arquitecturas y la verificación de especificaciones. Este enfoque respalda el requisito de escalabilidad de la solución, en la medida en que el modelo formal puede acompañar la evolución del sistema.

A diferencia de los trabajos citados, centrados en sistemas de tiempo real o en la infraestructura de herramientas, esta propuesta aplica VDM++ a la verificación de un algoritmo de planificación de rutas de transporte público, en el contexto específico de la ciudad de Arequipa.

\section{Requerimientos Funcionales}
\label{sec:requerimientos}

La propuesta plantea una plataforma web con dos funcionalidades principales: la búsqueda de rutas de transporte y el reporte colaborativo de incidencias urbanas. Los requerimientos funcionales que se describen a continuación constituyen la base de la especificación formal; los requerimientos RF1 y RF2 son los que serán especificados y verificados con VDM++ en el siguiente entregable.

\begin{table}[htbp]
\centering
\small
\caption{Requerimientos funcionales del sistema.}
\label{tab:rf}
\begin{tabular}{|c|p{6.6cm}|}
\hline
ID & Requerimiento funcional \\
\hline
RF1 & Búsqueda de rutas directas: para un origen y un destino dados, el sistema debe devolver las rutas cuyo recorrido pase a una distancia menor o igual a un radio definido (1000 m) de ambos puntos. \\
\hline
RF2 & Búsqueda de rutas combinadas: el sistema debe devolver combinaciones de dos rutas con un punto de transbordo viable, respetando una distancia máxima de transbordo (100 m) y una longitud mínima de segmento (500 m). \\
\hline
RF3 & Visualización en mapa: el sistema debe mostrar en un mapa interactivo el recorrido de las rutas y sus paraderos. \\
\hline
RF4 & Estimación de distancia y tiempo: el sistema debe calcular la distancia total del viaje y un tiempo estimado, incluyendo el tiempo de transbordo. \\
\hline
RF5 & Gestión de datos de rutas: el sistema debe permitir cargar rutas desde archivos KML y asociar los paraderos cercanos a cada recorrido. \\
\hline
\end{tabular}
\end{table}

Los requerimientos RF1 y RF2 se formalizarán en VDM++ definiendo los tipos de datos (puntos, recorridos), las invariantes del dominio (longitud mínima del recorrido, validez de coordenadas) y las operaciones de búsqueda con sus correspondientes pre y postcondiciones. Por ejemplo, la postcondición asociada a RF1 establece que toda ruta directa devuelta debe encontrarse a una distancia menor o igual al radio definido tanto del origen como del destino.

\FloatBarrier
\section{Resultados}
\label{sec:resultados}

% Sección correspondiente al Tercer Entregable.
% Aquí se presentarán la validación del modelo VDM++ (obligaciones de prueba y casos de prueba).

\endinput

\section{Discusión}
\label{sec:discusion}

% Sección correspondiente al Tercer Entregable.

\endinput

\section{Conclusiones}
\label{sec:conclusiones}

% Sección correspondiente al Tercer Entregable.

\endinput

\section{Limitaciones}
\label{sec:limitaciones}

% Sección correspondiente al Tercer Entregable.

\endinput

\section{Trabajos Futuros}
\label{sec:trabajos_futuros}

% Sección correspondiente al Tercer Entregable.

\endinput


%% References with BibTeX database:

\bibliographystyle{elsarticle-num}
\bibliography{references}

\end{document}
